\subsection{带算子的向量丛}

设 G 为李群，X 为 \(C^{r}\) 类流形，且 \((g, x) \mapsto gx\) 是 \(C^{r}\) 类的 G 在 X 上的左作用律。设 E 为 \(C^{r}\) 类向量丛，底空间为 X，且 \(\pi: E \to X\) 为 E 到 X 的投影。对于所有 \(x \in X\)，令 \(E_{x}\) 为 E 在 x 处的纤维。令 \((g, u) \mapsto gu\) 为 G 在 E 上的左作用律，使 \(\pi\) 与 G 在 X 和 E 上的作用相容。对于所有 \(g \in G\) 和 \(x \in X\)，将其在 \(E_{x}\) 上的限制 \(u \mapsto gu\) 记为双射 \(\psi_{g,x}\)，它从 \(E_{x}\) 到 \(E_{gx}\)。我们假定，对于所有 \(g \in G\) 和 \(x \in X\)，\(\psi_{g,x}\) 是连续线性的，因而是可赋范空间 \(E_{x}\) 到可赋范空间 \(E_{gx}\) 的同构。

令 \(\phi\) 为自同构 \((g,x)\mapsto (g,gx)\)（定义在流形 \(\mathrm{G}\times \mathrm{X}\) 上）。令 \(p\) 为 \(\mathrm{G}\times \mathrm{X}\) 到 \(\mathbf{X}\) 的典范投影，\(\mathrm{E}'\) 为 \(\mathrm{E}\) 在 \(p\) 下的逆像。令 \(\psi \colon \mathrm{E}'\to \mathrm{E}'\) 为各映射 \(\psi_{g,x}\colon \mathrm{E}_{(g,x)}'\to \mathrm{E}_{(g,gx)}'\) 的和。

定义 4. 如果 \(\psi\) 是向量丛的 \(\phi\)-态射且属于 \(\mathbf{C}^r\) 类，则 \(\mathbf{E}\) 称为向量 \(\mathbf{G}\)-丛，类为 \(\mathbf{C}^r\)。

换言之，E 是 \(C^{r}\) 类向量 G-丛，当且仅当对于所有 \((g_{0}, x_{0}) \in \mathbf{G} \times \mathbf{X}\)，以下条件成立：存在 \((g_{0}, x_{0})\) 在 \(G \times X\) 中的一个开邻域 U，使得，如果 \(E^{\prime} \mid U\)（相应地，\(E^{\prime} \mid \phi(U)\)）借助向量图识别为纤维 M（相应地，N）的平凡向量丛，则映射 \((g, x) \mapsto \psi_{g,x}\) 从 U 到 \(\mathcal{L}(M, N)\) 是 \(C^{r}\) 类的。

映射 \(\psi\) 显然是双射，由上述局部判据可知 \(\psi^{-1}\) 是向量丛的 \(\phi^{-1}\)-态射，因此 \(\psi\) 是向量丛的 \(\phi\)-同构。

以 X 为底空间的平凡向量 G-丛是向量丛 \(X \times F\)（其中 F 是完备可赋范空间），并带有作用律 \((g, (x, f)) \mapsto (gx, f)\) 在 G 的 \(X \times F\) 上的作用。

仍假定定义 4 前的假设和记号，并进一步取 \(\tau\) 为对同构的 \(C^r\) 类向量函子（Differentiable and Analytic Manifolds, R, 7.6.6）。则 \(\tau E\) 是底空间为 \(X\) 的向量丛。对于所有 \(x \in X\)，其纤维 \((\tau E)_x\) 等于 \(\tau (E_x)\)。对于所有可赋范空间 \(N_1, N_2\)，令 \(\operatorname{Isom}(N_1, N_2)\) 表示从 \(N_1\) 到 \(N_2\) 的同构集合。如果 \(g \in G\)，则

\[
\tau (\psi_ {g, x}) \in \operatorname{Isom} ((\tau E) _ {x}, (\tau E) _ {g x}).
\]

\(\tau(\psi_{g,x})\) 定义了左作用律 \((g,u)\mapsto gu\) G 在 \(\tau\) E 上，且 \(\tau\) E 到 X 的典范投影与 G 在 X 和 \(\tau\) E 上的作用相容。

命题 15. 如果 E 是 \(\mathbf{C}^r\) 类向量 G-丛，则 \(\tau \mathbf{E}\) 是 \(\mathbf{C}^r\) 类向量 G-丛。

设 \(g_0, x_0\)、U、M、N 如定义 4 后的段落所述。映射 \((g, x) \mapsto \tau(\psi_{g,x})\) 从 U 到 \(\mathcal{L}(\tau\mathrm{M}, \tau\mathrm{N})\)，是映射 \((g, x) \mapsto \psi_{g,x}\) 从 U 到 \(\mathcal{L}(\mathrm{M}, \mathrm{N})\) 与映射 \(f \mapsto \tau(f)\) 从 \(\operatorname{Isom}(\mathrm{M}, \mathrm{N})\) 到 \(\operatorname{Isom}(\tau\mathrm{M}, \tau\mathrm{N})\) 的复合；这两个映射都是 \(C^r\) 类的，因而复合也如此，从而得到命题。

命题 16. 设 G 为李群，X 为 \(\mathbf{C}^r\) 类流形（\(r \geqslant 2\)），且 \((g, x) \mapsto gx\) 是 G 在 X 上的 \(\mathbf{C}^r\) 类左作用律；因此通过传递结构，G 在 TX 上有一个左作用律。在此作用律下，TX 是 \(\mathbf{C}^{r-1}\) 类向量 G-丛。

令 \(\mathrm{pr}_1\)（相应地，\(\mathrm{pr}_2\)）为从 \(\mathrm{G} \times \mathrm{X}\) 到 \(\mathrm{G}\)（相应地，\(\mathrm{X}\)）的典范投影，并令 \(\mathrm{E}_1\)（相应地，\(\mathrm{E}_2\)）为关于 \(\mathrm{pr}_1\)（相应地，\(\mathrm{pr}_2\)）的 TG（相应地，TX）的逆像。于是向量丛 \(\mathrm{T}(\mathrm{G} \times \mathrm{X})\) 是 \(\mathrm{E}_1\) 和 \(\mathrm{E}_2\) 的直和。令 \(i: \mathrm{E}_2 \to \mathrm{T}(\mathrm{G} \times \mathrm{X})\) 和 \(q: \mathrm{T}(\mathrm{G} \times \mathrm{X}) \to \mathrm{E}_2\) 为由此直和分解定义的典范向量丛态射。令 \(\phi\) 为映射 \((g, x) \mapsto (g, gx)\)，它从 \(\mathrm{G} \times \mathrm{X}\) 到 \(\mathrm{G} \times \mathrm{X}\)。然而，定义 4 中记作 \(\psi\) 的映射（其中取 \(\mathrm{E} = \mathrm{TX}\)）正是 \(q \circ \mathrm{T}(\phi) \circ i\)。但是 \(\mathrm{T}(\phi)\) 是一个 \(\phi\)-态射，并且作为向量丛态射的类为 \(\mathrm{C}^{r-1}\)（Differentiable and Analytic Manifolds, R, 8.1.2）。

推论。若 \(\tau\) 是对同构的 \(\mathbf{C}^r\) 类向量函子，则 \(\tau(\mathrm{TX})\) 是类为 \(\mathbf{C}^{r-1}\) 的向量 G-丛。

这由命题 15 和 16 得出。

注 1. 若 \(\tau\) 是有限维情形下对同构的 \(C^r\) 类向量函子，且 \(E\) 具有有限秩，则 \(\tau E\) 的定义类似，命题 15 仍然成立；只要 \(X\) 是有限维的，命题 16 的推论仍然成立。

例子。取命题 16 的假设和记号，并令 F 为完备可赋范空间。则 \(\mathcal{L}((\mathrm{TX})^p; \mathrm{F})\) 是类为 \(\mathrm{C}^{r-1}\) 的向量 G-丛；如果 K 的特征为零或 X 是有限维的，则 \(\operatorname{Alt}^p(\mathrm{TX}; \mathrm{F})\) 也是如此（参见~Differentiable and Analytic Manifolds, R, 7.7, 7.8）。如果 X 是有限维的，则 \(\bigotimes^p(\mathrm{TX}) \otimes \bigotimes^q(\mathrm{TX})^*\) 是类为 \(\mathrm{C}^{r-1}\) 的向量 G-丛。

命题 17. 设 G 是李群，X 是 G 的左李齐性空间，\(x_0\) 是 X 中的一点，\(G_0\) 是 G 中 \(x_0\) 的稳定子，E 和 E' 是基空间为 X 且类为 \(C^{r}\) 的左向量 G-丛，\(E_0\)（相应地，\(E_0'\)）是 E（相应地，E'）在 \(x_0\) 处的纤维，且 f 是 \(\mathcal{L}(E_0, E_0')\) 中满足 \(f(gu) = g(f(u))\)（对所有 \(u \in E_0\) 和 \(g \in G_0\)）的元素。则存在唯一一个从 E 到 E'、与 G 的作用相容并延拓 f 的态射。这个态射的唯一性是显然的。下面证明其存在性。令 \(g\)、\(g'\) 为 G 中的元素，且 \(u \in E_0\) 满足 \(gu = g'u\)。则 \(g'^{-1}g \in G_0\) 且 \(g'^{-1}gu = u\)，从而 \(g'^{-1}gf(u) = f(u)\)，即 \(gf(u) = g'f(u)\)。因此定义映射 \(\phi\)，其从 E 到 \(E'\)，规定 \(\phi(gu) = gf(u)\)。显然，此映射延拓 \(f\)，并且与 G 的作用相容。我们证明 \(\phi\) 是类为 \(C^r\) 的向量丛态射。令 \(x_1 \in X\)。存在 X 中 \(x_1\) 的开邻域 V 以及 G 的子流形 W，使得映射 \(g \mapsto gx_0\) 是 W 到 V 的同构 \(\theta\)，类为 \(C^r\)。将 V 和 W 缩小后，可设：

\begin{enumerate}
\def\labelenumi{(\arabic{enumi})}
\item
  E \textbar{} V（相应地，E' \textbar{} V）被识别为纤维为 M（相应地，M'）的平凡向量丛；
\item
  若 \(\psi_{g}\)（相应地，\(\psi_{g}^{\prime}\)）表示映射 \(u\mapsto gu\)，其定义域为 \(\mathrm{E}_0\)（相应地，\(\mathrm{E}_0^{\prime}\)），值域为 \(\mathrm{E}_{gx_0}\)（相应地，\(\mathrm{E}_{gx_0}\)），则映射 \(g\mapsto \psi_g\) 和 \(g\mapsto \psi_g^{-1}\)（相应地，\(g\mapsto \psi_g^{\prime}\) 和 \(g\mapsto \psi_g^{\prime -1}\)）从 W 到 \(\mathcal{L}(\mathrm{E}_0,\mathrm{M})\) 和 \(\mathcal{L}(\mathrm{M},\mathrm{E}_0)\)（相应地，\(\mathcal{L}(\mathrm{E}_0',\mathrm{M}')\) 和 \(\mathcal{L}(\mathrm{M}',\mathrm{E}_0'))\)，都是类为 \(C^r\) 的。
\end{enumerate}

对于 \(x \in V\)，令 \(\phi_x: M \to M'\) 为 \(\phi\) 在 \(E_x = M\) 上的限制。于是 \(\phi_x\) 由以下映射复合而成：

\begin{enumerate}
\def\labelenumi{(\arabic{enumi})}
\item
  映射 \((\psi_{\theta^{-1}(x)})^{-1}\)，它从 M 到 \(E_{0}\)；
\item
  从 \(E_{0}\) 到 \(E_{0}'\) 的映射 f；
\item
  映射 \(\psi_{\theta^{-1}(x)}^{\prime}\)，它从 \(\mathbf{E}_0'\) 到 \(\mathbf{M}'\)。
\end{enumerate}

因此可见，映射 \(x \mapsto \phi_x\) 从 V 到 \(\mathcal{L}(\mathbf{M}, \mathbf{M}')\)，且是类为 \(\mathbf{C}^r\) 的。

推论 1. 设 \(\mathrm{E}_0^{\mathbb{G}_0}\) 是 \(\mathrm{E}_0\) 中在 \(\mathbf{G}_0\) 作用下不变的元素所成的集合。对所有 \(u \in \mathrm{E}_0^{\mathbb{G}_0}\)，令 \(\sigma_u\) 为从 X 到 E 的映射，由 \(\sigma_u(gx_0) = gu\)（对所有 \(g \in G\)）定义。

\begin{enumerate}
\def\labelenumi{(\roman{enumi})}
\item
  E 的 G-不变截面†都是类为 \(\mathbf{C}^r\) 的。
\item
  \(u \mapsto \sigma_u\) 是从 \(\mathrm{E}_0^{\mathrm{G}_0}\) 到 E 的 G-不变截面集合的双射。
\end{enumerate}

断言 (ii) 是显然的。为证明 (i)，只需证明每个截面 \(\sigma_{u}\) 都是类为 \(C^r\) 的。令 E' 为基空间 X、纤维为 \(\mathrm{E}_0^{\mathrm{G}_0}\) 的平凡 G-丛。令 \(f\) 为从 \(\mathrm{E}_0^{\mathrm{G}_0}\) 到 \(\mathrm{E}_0\) 的典范嵌入。根据命题 17，存在态射 \(\phi\)，它从 E' 到 E、与 G 的作用相容并延拓 \(f\)。若 \(u \in \mathrm{E}_0^{\mathrm{G}_0}\) 且 \(g \in \mathrm{G}\)，则

\[
\sigma_ {u} (g x _ {0}) = g u = g f (u) = \phi (g u) = \phi ((u, g x _ {0}))
\]

从而有 \(\sigma_u(x) = \phi((u, x))\)，对所有 \(x \in \mathbf{X}\)，这就证明了断言。

*例如，设 G 是有限维实李群，\(G_{0}\) 是 G 的紧李子群，X 是齐性空间 \(G/G_{0}\)。记 X 中 e 的典范像为 \(x_{0}\)。在 \(T_{x_{0}}(X)\) 上存在一个在 \(G_{0}\) 作用下不变的正定对称双线性型（Integration，第 VII 章，第 § 3 节，命题 1）。将上述结果应用于 \((\mathrm{TX})^{*}\otimes(\mathrm{TX})^{*}\)，可见 X 上存在一个在 \(G_{0}\) 作用下不变的解析黎曼度量。

推论 2. 假设 \(G_{0}\) 在 \(E_{0}\) 上平凡作用。令 \(E'\) 为基空间 X、纤维为 \(E_{0}\) 的平凡 G-丛。存在唯一一个从 E 到 \(E'\)、与 G 的作用相容并延拓 \(Id_{E_{0}}\) 的同构。

这立即由命题 17 得出。

注 2. 在本编号中，左作用律处处都可以用右作用律替代。

